\documentclass[../main.tex]{subfiles}
\begin{document}
%==========================================TIÊU ĐỀ TẬP========================================
\begin{table}[h]
    \begin{adjustwidth}{-1cm}{}
        \begin{tabular}{>{\centering\arraybackslash}p{6cm}|>{\raggedright\arraybackslash}p{12.5cm}}
            \multirow{5}{6cm}
            {\\[-40pt]
                \flaregothic\fontsize{20pt}{20pt}\selectfont \centering
                \color{eptype}HISTOIRE\color{black}\\
                \freshbot\fontsize{72pt}{72pt}\selectfont
                \color{epnum}27\\[6pt]
                \cafeteria\fontsize{20pt}{20pt}\selectfont\color{epnum}(D-27)\color{black}\\[10pt]
                \sen\fontsize{18pt}{18pt}\selectfont
                \color{epattr}Mini-histoire
                \color{black}
            }
             &
            {\fotskip\fontsize{18pt}{18pt}\selectfont \ruby{四}{よん}\ruby{色}{しき}のholoForce}
            \\[6pt]
             & {\cafeteria\fontsize{18pt}{18pt}\selectfont Four Colors of holoForce}      \\[4pt]
             & {\vnmsans\fontsize{18pt}{18pt}\selectfont Bốn sắc màu holoForce}           \\[8pt]
             & {\caviar{\fontsize{14pt}{14pt}\selectfont Nhân vật: Theo từng câu chuyện}} \\[8pt]
        \end{tabular}
    \end{adjustwidth}
\end{table}
%=========================================TRUYỆN CHÍNH========================================
\color{black}

\txtbx[2]{
    {\color{vol} \uline{Tóm tắt:}} Nối tiếp cuộc tái ngộ đầy xúc động tại kho hàng ngoại ô Kawachi, Kson lần lượt có những buổi gặp gỡ riêng tư với từng thành viên của Thế hệ thứ Tư. Một chuyến ghé thăm của Kanata làm sáng tỏ bí mật về bộ sưu tập merchandise và chiếc áo ``giẻ lau'' huyền thoại; lời cầu cứu dở khóc dở cười của nàng cựu Rồng khiến ``đùi cừu'' Watame bị xoay như chong chóng; đại tiệc takoyaki cay xè thực hiện trọn vẹn lời hứa năm xưa với công chúa Luna; và màn chăm sóc tận tình của ``ác ma'' Towa bị lật tẩy thành thiên sứ đích thực. Khép lại chuỗi ngày tái ngộ, cả năm thành viên holoForce cùng quây quần bên bếp than rực hồng trên sân thượng nhà Hikawa, sẻ chia những miếng thịt nướng giòn tan và ngắm nhìn muôn ngàn vì sao đêm.
}

\stpt{-Histoire 1-}
\stcap{Báu vật của Tổng trưởng}
\begin{center}
    {\caviar\fontsize{12pt}{12pt}\selectfont Nhân vật: \emp{kanata} \emp{kson} \emp{kuon2}}
\end{center}

\pov{Amane Kanata}


Suốt cả tuần lễ sau buổi chiều lao động náo loạn ở kho hàng ngoại ô Kawachi, đầu óc tôi vẫn không tài nào dứt ra khỏi một câu nói.

Đó là lời trần tình vô tư của Nene-chi trong bữa cơm hôm nọ:

\txtbx{``\dots Và đặc biệt là cực kỳ thích sưu tầm đồ lưu niệm có in hình của Kanata-senpai ạ!''}

Một con Rồng cao to, ăn to nói lớn, phong thái ngổ ngáo như đại tỷ giới giang hồ, câu cửa miệng lúc nào cũng là tiếng Anh chửi thề bạt mạng\dots\ lại đi cuồng sưu tầm merchandise của tôi á?!

Cái tên suốt ngày trêu tôi là ``khỉ đột lực điền'', chê tôi ``ngực lép như tấm ván phẳng lì'' ấy sao? Làm sao một kẻ chuyên đi cà khịa như cậu ta lại chịu bỏ tiền túi ra rước ba cái đồ lỉnh kỉnh in mặt tôi về nhà được chứ?

Không thể nào. Chắc chắn con bé Nene lại thêu dệt, hoặc hoa mắt nhìn gà hóa cuốc rồi.

Thế nhưng, sự tò mò cứ như đàn kiến bò râm ran trong dạ dày, thôi thúc tôi không yên. Đến trưa ngày thứ Bảy, viện cớ mang một túi bánh su kem matcha thơm phức sang cảm ơn vì bữa nước giải khát ở kho hôm nọ, tôi quyết định bắt chuyến tàu điện thẳng tiến về vùng ngoại ô Kawachi, ghé thăm căn nhà của gia đình Hikawa.

\ct

\begin{center}
    \uline{Ký túc xá nhà Hikawa - Thị trấn Kawachi.}
\end{center}

Cửa tiệm viễn thông ở tầng một mở rộng đón gió trưa. Anh Kuon đang ngồi sau quầy thanh toán, tay gõ bàn phím lạch cạch xử lý hợp đồng, trên bàn vẫn còn nguyên cốc cà phê đen bốc khói mỏng. Thấy tiếng chuông cửa leng keng, anh ngước lên, đôi lông mày hơi nhướn ra vẻ ngạc nhiên:

``Ủa, Kana-chan đấy à? Nay rảnh rỗi ghé qua đây có việc gì thế em?''

Tôi giơ giơ túi bánh su kem trong tay, cố nở một nụ cười thật tự nhiên: ``Em tiện đường ghé qua gửi ít bánh cho cả nhà thôi ạ. Mà\dots\ Kson-chan có nhà không anh?''

Khóe môi anh Kuon bất giác nhếch lên thành một nụ cười đầy ẩn ý. Anh khẽ thở dài một tiếng, lắc lắc đầu: ``Có, đang ở trên phòng tầng bảy đấy. Mà này\dots\ anh khuyên thật là chuẩn bị tâm lý trước khi bước vào nhé.''

Tôi tròn mắt: ``Dạ? Chuẩn bị tâm lý gì cơ anh?''

``Thì\dots\ cứ lên đó rồi khắc biết,'' anh nhún vai, quay lại màn hình máy tính, buông thõng một câu: ``Anh không chịu trách nhiệm nếu em bị sốc văn hóa đâu đấy.''

Lời cảnh báo mơ hồ của anh Tổng quản lý càng khiến tim tôi đập thình thịch. Tôi chào anh rồi rón rén bước lên những bậc cầu thang gỗ dẫn lên các tầng trên.

Lên đến hành lang tầng bảy, không gian tĩnh lặng lạ thường. Cánh cửa phòng ngoài cùng bên phải hơi khép hờ, để lộ một khe hẹp le lói ánh đèn vàng ấm áp. Tôi đứng trước cửa, giơ tay định gõ, nhưng rồi lại nảy ra ý định tinh nghịch muốn hù cho cô bạn một phen giật thót.

Tôi nhẹ nhàng đẩy mép cửa, bước lùi vào trong, hít sâu một hơi rồi định hô lớn: ``Good afternoon, motherfucker--''

Thế nhưng, âm thanh chưa kịp thoát ra khỏi lồng ngực thì cổ họng tôi đã nghẹn ứ lại.

Căn phòng vắng tanh. Kson không có ở đây, có lẽ cậu ấy đang dưới bếp chung nấu nước hoặc chạy ra cửa hàng tiện lợi gần nhà.

Nhưng đó không phải là lý do khiến tôi đứng chết trân như một pho tượng đá.

Mà là vì\dots\ những thứ đang hiện diện trước mắt tôi.

Trái ngược với phong cách áo da hầm hố, gậy bóng chày và huy hiệu kim loại gai góc đặc trưng của một cô nàng Tổng trưởng xã hội đen, toàn bộ góc tường bên trái và dãy kệ gỗ cạnh bàn làm việc của cậu ấy\dots\ trông hệt như một viện bảo tàng chuyên đề về Amane Kanata.

Tôi dụi mắt tới ba lần, cứ ngỡ mình hoa mắt vì tụt đường huyết.

Trên kệ, hàng chục chiếc standee acrylic của tôi được xếp theo thứ tự thời gian vô cùng ngăn nắp: từ bộ trang phục ra mắt ban đầu với chiếc áo choàng xanh lam cánh thiên sứ, cho đến bộ kimono đón năm mới, trang phục bơi 3D mùa hè, rồi cả bộ váy idol lấp lánh của buổi hòa nhạc lớn. Bên cạnh đó là cả một bộ sưu tập móc khóa, huy hiệu phiên bản giới hạn, khăn choàng cổ in hình linh vật thiên sứ, và ngay trên đầu giường ngủ là hai chú búp bê nhồi bông Upabikari tròn ủng đang nằm ngay ngắn cạnh chiếc gối ôm.

Chưa dừng lại ở đó, trên giá sách còn dựng trang trọng chiếc đĩa CD single có chữ ký in nổi của tôi, bên cạnh là một xấp bìa thẻ bo góc được bọc màng bảo vệ sáng loáng không tì vết.

Máu trong người tôi dồn hết lên mặt, nóng bừng như bị hơ trên lửa. Nene-chi không hề nói điêu một nửa lời!

Con Rồng này\dots\ Cậu ta âm thầm lùng sục, đặt mua hết sạch tất cả merchandise của tôi từ lúc nào không hay thế này?!

Tôi ôm lấy hai bên má đang đỏ gay gắt, vừa bàng hoàng vừa lúng túng không biết nên giấu mặt vào đâu. Thế nhưng, khi ánh mắt tôi vô tình lướt lên vị trí cao nhất trên kệ tường trung tâm -- nơi được bố trí dải đèn LED chiếu rọi trang trọng nhất -- toàn thân tôi bỗng khựng lại.

Ở đó, được bảo quản cẩn thận bên trong một khung kính acrylic trong suốt đóng nẹp gỗ sang trọng, là một chiếc áo phông trắng cộc tay.

Chất vải cotton hơi sờn nhẹ theo năm tháng. Và nổi bật ngay giữa ngực áo là những nét chữ viết tay bằng mực vẽ vải màu xanh tím nguệch ngoạc, nét đậm nét nhạt xiêu vẹo ``PP-Tenshi''.

Tôi nín thở. Đồng tử co thắt dữ dội.

Ký ức của mùa hè năm 2020 - cái ngày định mệnh ở văn phòng - lập tức dội về trong tâm trí tôi rõ mồn một như một thước phim quay chậm.

\ct

\begin{center}
    \uline{Hồi tưởng\\Mùa hè năm 2020 - Văn phòng \hll.}
\end{center}

Hôm ấy, tôi đã đạp tung cánh cửa văn phòng với tất cả niềm kiêu hãnh của một nhà sáng tạo trẻ tuổi. Vầng hào quang trên đầu quay tít, đôi cánh nhỏ phe phẩy, tôi giơ cao chiếc áo phông trắng mà mình đã thức trắng cả đêm cặm cụi vẽ từng nét mực đến cứng đơ cả khớp ngón tay:

``Em đã tự tay làm một sản phẩm mới! Lần này em đầu tư công sức lắm đấy!''

Chị Haato đã tấm tắc khen: ``Trông nó đáng yêu thật đấy nhỉ! Giống kiểu nghệ thuật đường phố ấy!''

Anh Kuon cũng gật gù động viên: ``Anh thấy sáng tạo đấy. Ít nhất thì em cũng không đổ bột mì lên người như lần trước.''

Thế nhưng, khi tôi tràn trề hy vọng quay sang Coco, con Rồng ấy chỉ khoanh hai tay trước ngực, nhìn chiếc áo bằng ánh mắt vô hồn của một kẻ đang ngắm nhìn bãi rác ven đường, rồi thản nhiên phán một câu xanh rờn làm tan nát cõi lòng tôi: ``Cái giẻ lau bẩn thỉu này là sao đây? Trông như ai đó đã giặt quần áo xong rồi dùng nó để lau sàn vậy.''

Tôi đã tức điên lên, nhảy cẫng lên máy tính gõ cụm từ ``giẻ rách'' (ボロ雑巾) trên thanh tìm kiếm PP Search hòng chứng minh con rồng mù thẩm mỹ. Nhưng hỡi ôi, thuật toán tàn nhẫn của internet lại trả về ngay kết quả đầu tiên chính là hình ảnh chiếc áo phông của tôi, kèm theo dòng mô tả: ``Áo thun giẻ lau - độc đáo, thời trang, và\dots\ lau sạch.''

Chị Haato ôm bụng nhịn cười đến run cả người. Anh Kuon xoa trán bất lực. Còn Coco thì nhấp một ngụm nước, tỉnh bơ chốt hạ: ``Ừ, đúng là đồ giẻ lau thật mà. Internet không bao giờ sai.''

Sau đó là một trận hỗn chiến long trời lở đất khi tôi nổi cơn thịnh nộ đuổi đánh con rồng chạy vòng quanh bàn làm việc. Chiếc áo phông bị tôi vứt lại chỏng chơ trên bàn trong nỗi phẫn uất tột cùng, và từ đó về sau, tôi chẳng buồn đụng tới nó nữa\dots

\ct

\begin{center}
    \uline{Hiện tại\\Phòng của Kson.}
\end{center}

Tôi run rẩy đưa bàn tay chạm nhẹ lên lớp kính acrylic mát lạnh.

Vết mực hơi nhòe ở góc dưới chân chữ ``P'' thứ hai do tôi vô tình quẹt tay vào hồi đó vẫn còn nguyên vẹn. Từng đường chỉ viền cổ áo hơi nhăn nheo đặc trưng\dots\

Chính là nó. Không thể lẫn vào đâu được.

Chiếc áo thun ``giẻ lau'' huyền thoại năm ấy.

Cái thứ mà con Rồng ngang tàng kia từng công khai chê bai thậm tệ trước mặt mọi người, ví von không khác gì mảnh giẻ rách dùng để chùi sàn nhà vệ sinh\dots\ vậy mà giờ đây, nó lại đang được lồng khung kính cẩn thận, ủi phẳng phiu từng nếp gấp và treo ở vị trí tôn nghiêm nhất trong căn phòng này?!

``Khoan đã\dots\ Con Thằn lằn bự con này\dots\ Rốt cuộc là thế nào?!'' tôi lẩm bẩm trong miệng, hai tai nóng bừng bừng.

*CẠCH!*

Tiếng tay nắm cửa vặn mở vang lên cắt ngang dòng suy nghĩ của tôi.

Kson bước vào, trên tay bưng một chiếc khay gỗ đựng hai cốc trà lúa mạch mát lạnh cùng một đĩa bánh gạo nướng giòn: ``Này Kanatan, Kuon-paisen bảo bà sang chơi nên tôi tiện thể xuống bếp lấy ít nước mát--''

Lời nói của cô nàng đột ngột đứt phéng giữa chừng.

Chiếc khay gỗ trên tay Kson chao đảo suýt rơi tuột xuống đất.

Ánh mắt hai đứa chúng tôi va vào nhau giữa không trung. Rồi từ ánh mắt của tôi, tầm nhìn của Kson chậm rãi trượt xuống theo hướng ngón tay trỏ của tôi đang chỉ thẳng vào khung kính trên tường.

Một giây trôi qua.

Hai giây trôi qua.

Ba giây tĩnh lặng như tờ, chỉ còn nghe thấy tiếng ve sầu râm ran ngoài cửa sổ buổi trưa.

Gương mặt thường ngày vốn ngầu lòi, phong trần, bất cần đời của cô nàng Tổng trưởng bắt đầu chuyển từ trắng bệch sang đỏ ửng, rồi bùng lên đỏ tía tai như một quả cà chua chín mọng.

``\eng{W-What the f*ck\dots?!}'' Kson buột miệng thốt lên một câu tiếng Anh, vội vã đặt mạnh khay nước xuống bàn làm nước trà sánh cả ra ngoài. Cậu ta lao vụt tới như một cơn gió, giang rộng hai cánh tay to lớn cố che kín lấy khung kính trên tường: ``K-Kanatan! Sao bà vào phòng người ta mà không gõ cửa hả?! Ai cho phép bà đi lục lọi lung tung thế này?!''

Tôi từ từ hạ ngón tay xuống, hai tay chuyển sang chống nạnh, khóe môi nhếch lên một nụ cười vừa ranh mãnh vừa tràn ngập sát khí: ``Ồ? Tôi lục lọi lung tung á? Tôi đường hoàng bước vào phòng bằng cửa chính nhé! Nhưng mà Tổng trưởng Kson vĩ đại ơi\dots\ Bà có thể giải thích cho tôi biết cái thứ treo sau lưng bà là cái của nợ gì không?''

Kson liếc mắt nhìn ra sau lưng, rồi quay lại nhìn tôi, nuốt nước bọt cái ực, cố gượng gạo cãi chày cãi cối: ``T-Thì là cái áo phông chứ cái gì! Phòng của tôi, tôi thích treo áo gì mà chẳng được!''

``Áo phông bình thường á?'' tôi bước lên một bước, ép sát khiến Kson bất giác lùi lại nửa nhịp. Tôi gằn từng chữ: ``Năm xưa ai là người mồm to chê bai rằng: \textit{`Cái giẻ lau bẩn thỉu này là sao đây? Trông như ai đó giặt quần áo xong rồi dùng nó để lau sàn vậy'} hả?! Ai là người hả hê tuyên bố \textit{`Internet không bao giờ sai'} khi thấy kết quả tìm kiếm giẻ rách hả?!''

Bị bóc mẽ không sót một chữ nào từ ký ức cũ, cô nàng cựu Rồng lắp bắp, hai má đỏ gay gắt: ``Thì\dots\ thì nó đúng là giẻ lau chứ sao! Nhìn cái nét chữ nguệch ngoạc của bà xem, có khác gì chữ gà bới không cơ chứ!''

``Giẻ lau mà bà đóng khung kính acrylic viền gỗ sang trọng thế kia à?!'' tôi chỉ thẳng vào lớp kính sáng bóng. ``Giẻ lau mà bà hút ẩm, ủi phẳng phiu, lau bụi không sót một hạt cát thế kia à?! Còn cả đống standee với búp bê Upabikari trên giường kia nữa, bà tính giải thích thế nào đây hả con Thằn lằn kia?!''

Kson bị dồn vào chân tường, bèn vò đầu bứt tai, mặt mũi méo xệch vì xấu hổ. Cuối cùng, cô thở hắt ra một hơi dài, buông thõng hai tay chịu trận: ``Thôi được rồi! Tôi nhận! Tôi thua bà rồi đấy, đồ Khỉ đột đanh đá!''

Tôi bật cười đắc thắng: ``Mau khai thật đi! Bà lén nhặt chiếc áo này về từ lúc nào?''

Kson kéo chiếc ghế xoay lại ngồi phịch xuống, hai tay ôm trán, giọng lí nhí chẳng khác nào một đứa trẻ vừa bị mẹ bắt quả tang ăn vụng kẹo: ``Thì\dots\ hôm đó sau khi rượt đuổi tôi xong, bà hậm hực đùng đùng bỏ về, để quên luôn chiếc áo trên bàn văn phòng chứ đâu. Lát sau tôi thấy Kuon-paisen cầm chiếc áo lên, lẩm bẩm bảo định đem đi\dots\ làm giẻ lau bảng kính văn phòng thật.''

Tôi trố mắt: ``Cái gì cơ?! Kuon-senpai dám lấy áo của tôi đi lau bảng á?!''

``Thì đấy!'' cậu ấy bĩu môi. ``Thấy anh ấy chuẩn bị nhúng nước xà phòng, tôi mới giật phắt lại áo từ tay. Dù gì cũng là công sức của bà thức trắng cả đêm vẽ đến cứng đơ cả tay\dots\ để anh ấy đem đi chùi mực viết bảng thì uổng quá. Thế là tôi lén nhét vào ba lô mang về nhà.''

Giọng nói của cô nàng nhỏ dần lại, mang theo một chút ngượng ngùng chân thành: ``Về đến nhà, tôi giặt sạch sẽ rồi gấp gọn cất dưới đáy tủ. Sau này dọn qua bao nhiêu chỗ ở, rồi cả lúc\dots\ lúc tôi tốt nghiệp rời đi, chiếc áo đó vẫn luôn nằm trong vali của tôi. Đến khi chuyển về ký túc xá nhà Hikawa này, tôi thấy cái tường trống trải quá, với lại\dots\ cũng nhớ mấy trò quậy phá ngày xưa của bọn mình, nên mới đem ra đóng khung treo lên. Chỉ có thế thôi! Hết chuyện rồi đấy!''

Căn phòng bỗng chốc chìm vào một khoảng lặng êm đềm.

Tôi đứng nhìn người bạn trước mặt. Dưới mái tóc nhuộm highlight cá tính và bộ trang phục giang hồ bụi bặm, người đang ngồi kia vẫn vẹn nguyên là con Rồng ngốc nghếch, vụng về và giàu tình cảm của ngày nào.

Cái con người ngoài mặt lúc nào cũng bộc trực, ăn nói bạt mạng và chuyên gia châm chọc người khác, nhưng trong sâu thẳm tâm can lại luôn nâng niu từng kỷ niệm vụn vặt nhất của chúng tôi suốt ngần ấy năm trời chia cách.

Một làn sóng ấm áp dâng trào trong lồng ngực, khiến sống mũi tôi cay cay.

Tôi bước lại gần chiếc bàn, nhẹ nhàng rút một chiếc bút dạ lông dầu màu đen từ trong túi xách ra.

Kson ngước lên, chớp mắt nghi hoặc: ``Bà định làm gì đấy?''

``Mở khung kính ra,'' tôi mỉm cười, giọng dịu dàng.

``Hả? Để làm gì cơ?''

``Thì ký tên chứ làm gì!'' tôi hất cằm, làm bộ kiêu hãnh. ``Hàng thủ công độc bản giới hạn duy nhất một chiếc trên toàn thế giới của tác giả Amane Kanata cơ mà! Không có chữ ký xác thực của tác giả thì làm sao gọi là báu vật vô giá được chứ?''

Kson ngơ ngác nhìn tôi một hồi, rồi trên khóe môi cô nàng bất chợt nở nụ cười rạng rỡ quen thuộc -nụ cười toe toét ngập tràn ánh nắng của một kẻ ngang tàng:

``Hừ, bà chỉ khéo vẽ chuyện thôi\dots''

Nói thì càu nhàu thế, nhưng đôi tay lực lưỡng của Kson lại cẩn thận gỡ chốt khung kính xuống, mở toang mặt kính để lộ lớp vải áo mềm mại bên trong.

Tôi cúi người, nắn nót đặt bút ký một nét chữ thật đẹp kèm theo hình vẽ một thiên sứ nhỏ có đôi cánh tròn xoe ngay dưới dòng chữ ``PP-Tenshi''. Và ở góc dưới cùng, tôi hí hửng viết thêm một dòng chữ nhỏ:

\begin{center}
    \textit{``Tặng cho con Thằn lằn ngốc chuyên sưu tầm giẻ lau - Mãi mãi là bạn thân nhất!''}
\end{center}

Kson nhìn thấy dòng chữ ấy liền ôm trán kêu trời: ``Này nhé! Ai cho phép bà bôi bác tôi lên tác phẩm nghệ thuật hả?!''

``Kệ tôi! Bút trong tay tôi thì tôi viết!'' tôi lè lưỡi trêu chọc.

Hai chúng tôi nhìn nhau, rồi đồng thanh bật cười lớn. Tiếng cười giòn tan vang vọng khắp căn phòng nhỏ trên tầng bảy, xua tan đi tất cả những khoảng cách thời gian và không gian từng ngăn cách chúng tôi suốt một năm qua.

Kson cẩn thận đóng nẹp khung kính lại rồi treo nó trở về vị trí trung tâm. Sau đó, cô nàng quay sang vỗ vai tôi bôm bốp: ``Được rồi, coi như bà có lòng! Mau ngồi xuống đây xơi bánh su kem với uống trà mát đi, Đội trưởng Khỉ đột!''

``Này! Đã bảo là đừng có gọi tôi bằng cái biệt danh đó nữa cơ mà!''

Chúng tôi ngồi bệt xuống sàn gỗ, cùng nhau thưởng thức vị ngọt ngào của bánh ngọt và làn nước trà mát lạnh, rôm rả kể cho nhau nghe đủ thứ chuyện trên trời dưới biển.

Và trên bức tường cao phía sau, chiếc áo thun ``giẻ lau'' huyền thoại vẫn lặng lẽ tỏa sáng dưới ánh đèn, như một minh chứng bất diệt cho một tình bạn kỳ lạ nhưng vĩnh cửu.

\newpage

\stpt{-Histoire 2-}
\stcap{Đùi cừu bị nhờ vả}
\begin{center}
    {\caviar\fontsize{12pt}{12pt}\selectfont Nhân vật: \emp{watame} \emp{kson}}
\end{center}

\pov{Tsunomaki Watame}

Tôi đang ngồi nhâm nhi cốc trà ấm và thảnh thơi gặm mấy cọng khoai tây chiên giòn tan thì chiếc điện thoại trên bàn bỗng rung lên bần bật.

Màn hình nhấp nháy dòng tên hiển thị: ``Kson-chan''.

Tôi chớp chớp mắt ngạc nhiên. Kể từ sau buổi gặp ở kho hàng tuần trước, đây là lần đầu tiên cô nàng cựu Rồng chủ động liên lạc riêng với tôi. Tôi vội vàng nhấn nút nghe máy: ``A-lô, Coco-chi\dots\ à nhầm, Kson-chan đấy à?''

Đầu dây bên kia lập tức vang lên một tiếng rên rỉ đầy thảm thiết, nghẹn ngào như thể chủ nhân của nó đang ở bên bờ vực sinh tử: ``\textbf{Watame\dots! Đùi Cừu ơi\dots! Cứu nguy cho tôi với\dots! Vụ này mà bà không ra tay thì tôi toi đời mất thôi\dots!}''

Tôi giật bắn mình, suýt nữa làm rơi cả cốc trà: ``Ể?! C-Có chuyện gì thế?! Đừng bảo là bà bị băng đảng nào đến trả thù hay là\dots\ là đói bụng quá rồi định bắt tôi làm món súp cừu hầm đấy nhé?!''

``Bậy bạ nào! Ai thèm ăn thịt bà lúc này chứ!'' giọng Kson gắt lên giữa những tiếng thở dốc. ``Nói chung là khẩn cấp lắm rồi! Mau sang phòng tôi ở tầng bảy nhà Hikawa ngay, nhớ mang theo cả sổ tay công thức nấu ăn của bà nữa! Nhanh lên kẻo không kịp!''

*TÚT\dots\ TÚT\dots\ TÚT\dots*

Cuộc gọi bị ngắt cái rụp. Tôi ngơ ngác nhìn màn hình điện thoại, lòng dạ bồn chồn không yên. Mang theo sổ công thức nấu ăn á? Không lẽ có chuyện gì liên quan đến bếp núc sao?

Nghĩ đến cảnh người bạn to xác đang quằn quại trong đau đớn, tôi không dám chần chừ thêm một giây nào. Tôi vội vã nhét cuốn sổ tay vào túi xách, khoác chiếc áo choàng len quen thuộc rồi hớt hải bắt chuyến tàu điện sớm nhất hướng về Kawachi.

\ct

\begin{center}
    \uline{Ký túc xá nhà Hikawa - Phòng của Kson.}
\end{center}

Tôi thở hổn hển xô cửa bước vào, hét lớn: ``Kson-chan! Bà có sao không?! Tôi mang thuốc men với--''

Lời nói nghẹn lại ở cổ họng khi tôi nhìn thấy cảnh tượng trước mắt.

Cô nàng Tổng trưởng cao lớn đang nằm bò soài trên bàn làm việc, hai tay ôm chặt lấy bụng, mặt mũi nhăn nhó như một quả mướp đắng. Xung quanh bàn là một ``nghĩa địa'' rác thải ẩm thực: ba hộp bánh pizza đông lạnh đã cạn kiệt, vỏ mì ăn liền chất thành chồng cao ngất ngưởng, và không dưới năm lon nước ngọt có gas rỗng không nằm lăn lóc dưới sàn.

Tôi trố mắt kinh ngạc: ``Trời đất ơi\dots\ Chuyện gì đã xảy ra ở đây thế này?!''

Kson ngước đôi mắt thâm quầng lên nhìn tôi, giọng nói yếu ớt như sắp tắt thở: ``Đùi Cừu\dots\ bà đến rồi đấy à\dots\ Cứu tôi\dots\ Bụng tôi nó đang biểu tình dữ dội quá\dots''

Sau một hồi gặng hỏi, tôi mới vỡ lẽ ra đầu đuôi câu chuyện dở khóc dở cười.

Hóa ra từ ngày dọn về đây ở tự lập và bận rộn với công việc kho bãi, Kson hoàn toàn lười vào bếp. Suốt hai tuần liền, cậu ta chỉ sống dựa vào đồ ăn nhanh kiểu Mỹ, đồ đông lạnh mua sẵn và mì gói cay nồng. Kết quả là dạ dày của con Rồng bị viêm loét, người ngợm nóng trong phát sốt, da dẻ nổi mụn tùm lum. Sáng nay, anh Kuon và Suzuno lên kiểm tra phòng, thấy bãi chiến trường rác rưởi liền nổi trận lôi đình, anh Tổng liền ra tối hậu thư: nếu trong hôm nay Kson không chịu nấu ăn tử tế và tống khứ đống đồ ăn rác rưởi này đi, anh sẽ tịch thu toàn bộ máy tính lẫn TV chơi game!

Bị dồn vào bước đường cùng, nàng cựu Rồng chỉ còn biết cầu cứu ``đầu bếp số một của Gen 4'' là tôi.

Nghe xong, tôi vừa giận vừa buồn cười, bèn khoanh tay chống nạnh: ``Đấy! Ai bảo bà ăn uống vô tội vạ cho lắm vào! Người ngợm to xác thế kia mà ăn uống chẳng khác gì con nít cả!''

Kson chắp hai tay trước ngực, giương đôi mắt đáng thương nhất trần đời lên nhìn tôi: ``Tôi biết lỗi rồi mà\dots\ Bà dạy tôi nấu mấy món gì thanh đạm, dễ làm mà ngon lành đi! Nhìn mấy cọng rau xanh mà Suzu-chan nấu, tôi nuốt không vô\dots''

Nói đoạn, ánh mắt cô nàng bỗng trượt xuống cặp đùi tròn trịa mềm mại của tôi, khóe môi khẽ nhếch lên một nụ cười ranh mãnh cố hữu: ``Mà công nhận\dots\ nhìn cặp đùi cừu múp míp của bà, tự dưng tôi lại thấy thèm món thịt cừu nướng sốt tương quá ta\dots''

Tôi lập tức nhảy lùi lại một bước, hai tay ôm chặt lấy chân, phồng má giận dỗi: ``\textbf{Này nhé! Watame hơm có nhỗi, và Watame hơm phải đồ ăn dự trữ đâu đấy nhá!}''

Kson ôm bụng bật cười khùng khục: ``Đùa tí thôi mà làm gì căng thế! Nào, sư phụ Watame, mau cứu cái dạ dày của tôi đi!''

\ct

Thế là hai chúng tôi cùng nhau kéo xuống gian bếp chung ở tầng sáu.

Tôi quyết định dạy cho Kson hai món ăn kinh điển của gia đình Nhật Bản: canh rau củ hầm thanh đạm \textit{Kenchinjiru} và món thịt bò xào hành tây sốt dashi ngọt dịu. Là \emph{thịt bò}, chứ không phải là \emph{thịt cừu} đâu.

``Nào, bước đầu tiên là phải sơ chế củ cải, cà rốt và khoai tây,'' tôi vừa nói vừa đưa cho Kson một con dao bếp và chiếc thớt gỗ. ``Bà thái hình quân cờ vừa ăn nhé!''

Kson nhận lấy con dao, tư thế cầm dao trông ngầu lòi chẳng khác nào một kiếm sĩ chuẩn bị chém gục đối thủ trên chiến trường. Rồi cô nàng vung tay chém:

*RẬP! RẬP! RẬP!*

Tiếng dao nện xuống thớt nghe chan chát như tiếng búa tạ gõ đinh. Từng khúc cà rốt bị chém văng tung tóe khắp sàn bếp, miếng thì to như nắm tay, miếng thì nát bấy như cám.

Tôi sợ thót tim, vội vã lao vào giữ chặt tay cô bạn to xác: ``Trời ơi là trời! Bà đang nấu ăn hay đang đi thanh toán môn hộ thế hả?! Dao dùng để cắt gọt nhẹ nhàng chứ có phải đao chém củi đâu! Nhìn tôi làm mẫu này!''

Tôi nhẹ nhàng cầm lấy bàn tay to lớn vụng về của Kson, kiên nhẫn hướng dẫn từng động tác đưa dao mềm mại, cách gọt vỏ củ sen và ngâm nước để không bị thâm. Cô bạn tròn mắt nhìn theo từng ngón tay thoăn thoắt của tôi, ngoan ngoãn gật đầu lia lịa như một đứa trẻ lần đầu đi học mẫu giáo.

Tiếp đó, tôi chỉ cho cô nàng cách đun nước dùng dashi từ tảo bẹ kombu và nấm đông cô, cách nêm nếm mirin, nước tương shoyu và một chút rượu nấu ăn để giữ được vị ngọt thanh nguyên bản của rau củ mà không cần dùng đến bột ngọt hay dầu mỡ công nghiệp.

Mùi thơm dịu dàng, ấm áp của nước dùng nhanh chóng lan tỏa khắp gian bếp, xua tan đi sự ngột ngạt của buổi trưa hè.

Sau gần một tiếng đồng hồ đánh vật với xoong nồi, mâm cơm thanh đạm cuối cùng cũng hoàn thành.

Hai bát canh \textit{Kenchinjiru} nghi ngút khói với đủ sắc màu của củ cải trắng ngần, cà rốt đỏ au, nấm nâu và đậu phụ mềm mượt, bên cạnh là đĩa thịt bò xào hành tây thơm lừng rắc thêm chút hạt mè rang vàng.

Kson nuốt nước bọt ừng ực. Cô nàng run run cầm thìa múc một ngụm canh đưa lên miệng thổi nhẹ rồi nhấp thử.

Một giây sau, đôi mắt đỏ rực của nàng cựu Rồng bỗng sáng bừng lên như hai ngọn đuốc:

``\textbf{\eng{Holy shit\dots!} Ngon khủng khiếp luôn Đùi Cừu ơi!}''

Cô nàng húp sùm sụp hết thìa này đến thìa khác, rồi gắp một miếng thịt bò xào cho vào miệng nhai ngấu nghiến. Vị ngọt thanh tao của rau củ hòa quyện với nước dùng ấm nóng như xoa dịu từng tế bào trong dạ dày đang rệu rã của cô. Gương mặt nhăn nhó lúc nãy hoàn toàn biến mất, thay vào đó là một vẻ mặt thỏa mãn ngập tràn hạnh phúc:

``Ấm bụng quá\dots\ Cảm giác như toàn bộ độc tố trong người vừa bị cuốn phăng đi hết ấy!''

Tôi mỉm cười, tự hào ưỡn ngực: ``Thấy chưa! Tôi đã bảo đồ ăn tươi ngon tự nấu là thần dược tốt nhất mà lại!''

Kson đặt bát xuống, vỗ mạnh vào vai tôi một cái suýt làm tôi sặc nước trà: ``Cảm ơn bà nhiều lắm nhé, Watame! Không có bà chắc hôm nay tôi bị Kuon-paisen tống cổ ra đường rồi. Đúng là Đùi Cừu số một của lòng tôi!''

Tôi cười híp mí, dù cái biệt danh ``Đùi Cừu'' nghe vẫn hơi nhột một chút, nhưng nhìn người bạn thân thiết lấy lại được nụ cười rạng rỡ và sức khỏe, trong lòng tôi dâng lên một niềm vui sướng khôn tả.

Quả nhiên, dù con Rồng ấy có đổi tên hay mang vẻ ngoài dữ dằn đến đâu, sự ấm áp giữa chúng tôi vẫn vẹn nguyên như những ngày đầu tiên.

\newpage

\stpt{-Histoire 3-}
\stcap{Đại tiệc Takoyaki}
\begin{center}
    {\caviar\fontsize{12pt}{12pt}\selectfont Nhân vật: \emp{luna} \emp{kson}}
\end{center}

\pov{Himemori Luna}

Hôm nay là một ngày vô cùng trọng đại đối với công chúa xứ Kẹo tui-nanora!

Cách đây gần một năm, trên bờ biển lộng gió trước giờ chia tay, người bạn Rồng Coco-chi đã dõng dạc hứa trước mặt tui rằng: \textit{``Bữa nào tôi sẽ mời bà ăn tiệc takoyaki ngoài không gian luôn!''}.

Và hôm nay, lời hứa hẹn vĩ đại ấy cuối cùng cũng đến ngày thực hiện-nora!

Tui diện bộ váy bèo nhún màu hồng kẹo ngọt lộng lẫy nhất, đội chiếc vương miện nhỏ xíu trên đầu, tay cầm cây quyền trượng kẹo mút tung tăng bước lên tầng bảy của ký túc xá nhà Hikawa. Với tư cách là một công chúa quý tộc, tui đến đây để kiểm tra tay nghề ẩm thực của kẻ hầu cận Rồng Kson-nanora!

*CẠCH!*

Tui đẩy cửa bước vào, cất giọng sang sảng đầy quyền uy: ``Công chúa Luna ngự giá rùi đây-nanora! Kson đâu mau ra nghênh đón-nora!''

Ngay giữa phòng khách, một chiếc máy nướng takoyaki điện cỡ lớn sáng loáng đặt chình ình trên bàn trà. Xung quanh bày la liệt các đĩa nguyên liệu tươi ngon: từng miếng râu bạch tuộc to sụ luộc chín đỏ au, bát bột bánh pha sẵn mịn màng, gừng đỏ băm nhỏ, hành hoa thái nhuyễn, và đặc biệt là một đĩa phô mai mozzarella béo ngậy được chất cao như ngọn núi nhỏ.

Kson đang đứng đó, đầu quấn dải băng vải đen, hai tay cầm hai que xiên tre nhọn hoắt, cười toe toét để lộ hàm răng trắng bóng: ``Chào mừng Công chúa Luna-tan đến với Bếp nướng Rồng Vũ Trụ! Hôm nay tôi sẽ đích thân phục vụ bà một bữa tiệc takoyaki đỉnh cao nhất thiên hà!''

Tui hài lòng gật đầu, nhảy tót lên chiếc ghế bành êm ái nhất trong phòng, bắt chéo hai chân rồi vung quyền trượng kẹo mút chỉ đạo: ``Tốt lắm-nora! Phải làm cho thựt ngon đó nha! Nếu bánh bị cháy khét hay ngụi ngắt là tui phạt đánh đòn đó-nanora!''

``Tuân lệnh công chúa!'' Kson cười khì, lập tức xắn tay áo vào việc.

*XÈOOOO\dots!*

Mặt chảo gang nóng rực reo lên vui tai khi Kson thoăn thoắt quét một lớp dầu ăn bóng loáng. Cô nàng đổ từng muôi bột bánh sánh mịn vào các khuôn tròn, rồi nhanh tay thả từng miếng bạch tuộc to đùng cùng gừng đỏ và hành hoa vào giữa.

Mùi thơm béo ngậy của bột và hải sản lập tức bốc lên ngùn ngụt, khiến bụng tui sôi lên ùng ục-nanora.

Hai chiếc que xiên tre trong tay Kson múa lượn thoăn thoắt như làm xiếc. Cậu ấy khéo léo lật từng viên bánh, vo tròn chúng lại thành những quả cầu vàng ươm, vỏ ngoài xém cạnh giòn rụm một cách hoàn hảo.

Thế nhưng, trong lúc tui đang mải mê ngắm nhìn kỹ nghệ điêu luyện ấy, con Rồng tinh quái kia bỗng lén lút rút từ túi quần ra một chai thủy tinh nhỏ có in hình đầu lâu xương chéo màu đỏ thẫm.

*TÁCH! TÁCH!*

Cậu ta lén nhỏ vài giọt chất lỏng màu đỏ vào hai viên bánh ở góc khuôn, miệng tủm tỉm cười đầy tà ác mà cứ ngỡ tui không nhìn thấy.

Chỉ vài phút sau, mẻ bánh đầu tiên đã hoàn thành. Kson gắp sáu viên takoyaki nóng hổi ra một chiếc đĩa thuyền gỗ xinh xắn, rưới đẫm sốt ngọt nâu óng ánh và xịt một đường mayonnaise trắng muốt lên trên, rồi cung kính dâng lên trước mặt tui: ``Mời Công chúa thưởng thức viên takoyaki vũ trụ đặc biệt nhất!''

Tui chớp mắt nghi hoặc, nhìn viên bánh vàng rộm thơm lừng đang bốc khói nghi ngút. Trông hấp dẫn quá chừng, chẳng có vẻ gì là nguy hiểm cả-nora.

Không thể cưỡng lại sự cám dỗ ngọt ngào, tui liền cầm xiên tre gắp lấy viên bánh mà Kson vừa nhiệt tình chỉ trỏ, thổi phù phù hai cái rồi há to miệng cắn ngập một nửa viên bánh.

Vỏ bánh giòn tan vỡ ra trong khoang miệng. Vị béo ngọt của sốt và dai giòn của bạch tuộc vừa chạm vào đầu lưỡi thì\dots

*BÙÙÙÙM!*

Một ngọn lửa vô hình bùng nổ dữ dội ngay giữa cổ họng tui!

Cái cay xé họng, cay buốt tận óc, cay như thể có một đàn rồng phun lửa trực tiếp vào trong dạ dày tui vậy!

Mặt tui trong tích tắc biến thành màu đỏ gay gắt như quả cà chua luộc chín. Khói dường như xì ra từ hai bên tai, hai hàng nước mắt lập tức tuôn rơi lã chã như mưa rào.

``\textbf{Oaaa! Cay quá-nanoraaaa! Cay cháy lữi rùi-nanora! Khụ khụ\dots!}'' tui nhảy dựng lên khỏi ghế bành, hai tay ôm chặt lấy miệng, giậm chân thình thịch xuống sàn nhà: ``\textbf{Đây là bom nhịt hạch chứ bánh zì-nora! Coco-chi tính mưu sát công chúa hả-nanoraaaa?!}''

Kson đứng cạnh đó ôm bụng cười ngặt nghẽo, cười đến nỗi chảy cả nước mắt: ``Hahaha! Trúng bẫy rồi! Đó là sốt ớt Ghost Pepper siêu cay nhập khẩu từ Mỹ đấy bà ơi! Hahaha!''

Nhìn tui vừa khóc nức nở vừa thở phì phò như con cá mắc cạn, cô nàng mới giật mình hốt hoảng nhận ra mình đùa hơi quá trớn. Kson cuống cuồng vứt que xiên sang một bên, lao vội ra tủ lạnh mở toang cửa: ``Ấy chết! Tôi xin lỗi! Đừng khóc nữa mà công chúa ơi! Uống sữa đi này! Sữa tươi lạnh đây!''

Cô nàng hớt hải rót đầy một cốc sữa đặc có đường ướp lạnh, rồi còn múc thêm một muỗng kem vani to tướng đút thẳng vào miệng tui.

Làn sữa ngọt lịm và vị mát lạnh của kem từ từ dập tắt ngọn lửa địa ngục trong miệng tui. Tui ngồi thụp xuống ghế, sụt sịt mũi, hai mắt đỏ hoe lườm con Rồng xấu tính một cái cháy mặt: ``Đồ tồi-nanora\dots\ Dám đầu độc công chúa\dots\ Tui méch anh Kuon cho bà coi-nora\dots''

Kson quỳ rạp một chân xuống sàn, chắp hai tay lại rối rít xin lỗi: ``Tôi thề danh dự là tôi không dám đùa nữa đâu! Lần này tôi làm bánh ngọt ngào chuẩn vị công chúa một trăm phần trăm cho bà luôn! Nín đi mà, thương thương!''

Thấy điệu bộ sợ sệt hối lỗi cuống cuồng của cô nàng to xác, tui mới nguôi ngoai phần nào, bĩu môi ra lệnh: ``Làm lại ngay-nanora! Cho thật nhiều phô mai vào đó-nora!''

Kson thở phào nhẹ nhõm, lập tức quay lại bếp nướng mẻ bánh thứ hai.

Lần này, cô nàng cẩn thận gấp mười lần: từng viên bánh được nhồi ngập tràn phô mai mozzarella kéo sợi, nướng vừa chín tới mềm mịn tan chảy, bên trên rưới nước sốt ngọt dịu thơm lừng và rắc một lớp cá bào bonito mỏng dính đang nhảy múa rì rào theo làn hơi nóng.

Tui dè dặt cắn thử một miếng nhỏ.

Vị béo ngậy ngọt ngào của phô mai quyện với vị đậm đà của sốt và thịt bạch tuộc giòn sần sật lan tỏa khắp khoang miệng. Ngon đến mức cả người tui như tan chảy ra-nanora!

Đôi mắt tui sáng lấp lánh, miệng nhai nhồm nhoàm phồng cả hai má: ``Ưm\dots! Ngon tuyệt cú mèo-nanora! Bánh giòn rụm bên ngoài mà bên trong béo ngậy lun-nora!''

Kson mỉm cười nhẹ nhõm, lấy một tờ khăn giấy nhẹ nhàng lau đi vệt sốt nâu dính trên khóe môi tui, giọng đầy cưng chiều: ``Ngon thì ăn nhiều vào, ăn cho hết cả đĩa này luôn đi đồ công chúa mít ướt!''

Tui cười tít mắt, vung cây kẹo mút lên cao đầy đắc thắng: ``Được rồi-nora! Đại tiệc Takoyaki ngoài vũ trụ thành công mỹ mãn! Tha tội chết cho nhà ngưi đó-nora!''

Căn phòng ngập tràn hương thơm của bánh nướng và tiếng cười trong trẻo của hai đứa. Lời hứa năm xưa tưởng chừng như xa vời giữa muôn trùng sóng gió, hóa ra lại được hoàn thành trọn vẹn và ngọt ngào đến thế.

\newpage

\stpt{-Histoire 4-}
\stcap{Thiên sứ khoác áo Ác ma}
\begin{center}
    {\caviar\fontsize{12pt}{12pt}\selectfont Nhân vật: \emp{towa} \emp{kson}}
\end{center}

\pov{Tokoyami Towa}

Tôi là Tokoyami Towa. Một Ác ma cao quý và tàn nhẫn đến từ Ma giới.

Đó là chân lý bất di bất dịch mà tôi luôn khắc ghi trong tim. Một ác ma thì phải gieo rắc nỗi sợ hãi, phải biết trừng phạt những kẻ lười biếng, vô kỷ luật và dám làm xáo trộn trật tự thế giới.

Và hôm nay, đối tượng cần phải nhận sự trừng phạt thích đáng từ tôi chính là: Kiryu Coco-- à không, bây giờ là Tổng trưởng Kson!

Chuyện là sáng nay tạt qua văn phòng, tôi tình cờ nghe anh Kuon than thở rằng Kson dạo này vừa làm việc kho bãi nặng nhọc, vừa cày cuốc thâu đêm suốt sáng đến mức đau nhức khắp người, lưng còng vai mỏi, phòng ốc thì bừa bộn như bãi rác mà chẳng chịu chăm sóc bản thân.

Một lối sống bệ rạc và vô kỷ luật như thế là hoàn toàn không thể chấp nhận được!

Tôi lập tức xách theo chiếc giỏ mây to đùng đựng đầy ``công cụ trừng phạt'', đằng đằng sát khí bước lên tầng bảy của tòa nhà Hikawa. Tôi sẽ cho con Rồng ngốc nghếch ấy biết thế nào là sự tàn bạo của một Ác ma chân chính!

*CẠCH!*

Tôi đẩy cửa bước vào, chống hai tay lên hông, hất cằm cất giọng lạnh lùng nhất có thể: ``Kson! Mau đứng nghiêm chỉnh lại cho tôi! Hôm nay tôi đến đây là để gieo rắc sự trừng phạt khủng khiếp nhất lên đầu bà đấy!''

Thế nhưng, đáp lại lời đe dọa đầy quyền uy của tôi chỉ là một tiếng rên rỉ yếu ớt: ``Ư\dots\ TMT đấy à\dots\ Cửa không khóa đâu, vào đi\dots\ Ôi cái lưng của tôi gãy làm đôi rồi\dots''

Kson đang nằm gục trên chiếc ghế xoay, một tay xoa xoa bả vai cứng đờ, mặt mũi bơ phờ vì thiếu ngủ. Xung quanh bàn làm việc là cả một mớ hỗn độn: tài liệu vứt lung tung, đệm lót ghế xô lệch, và chăn màn trên giường thì vo tròn lại như một tổ chim sẻ.

Nhìn cảnh tượng bê tha ấy, ngọn lửa phẫn nộ trong lòng tôi bùng lên dữ dội. Tôi giậm mạnh chân: ``Đã bảo là đừng có gọi tôi là TMT cơ mà! Và nhìn cái phòng này xem, bà có còn là con người nữa không hả?!''

``Thì tôi là Rồng mà\dots'' Kson lẩm bẩm chống chế.

``\textbf{Im miệng!} Sự trừng phạt của Ác ma bắt đầu từ bây giờ! Run sợ đi!''

Tôi lập tức xắn tay áo lên, bắt đầu thực hiện kế hoạch ``hành hạ'' tàn nhẫn của mình.

Bước thứ nhất: tôi lao vào dọn dẹp căn phòng bừa bãi với tốc độ của một cơn lốc xoáy. Tôi gom sạch từng vỏ chai nước bỏ vào túi rác phân loại, quét sạch bóng sàn nhà, gấp chiếc chăn bông vuông vắn như một chiếc bánh chưng rồi xếp lại gọn gàng trên đầu giường.

Tôi vừa lau sạch bụi bặm trên bàn làm việc vừa gắt gỏng: ``Tôi dọn sạch phòng cho bà là để bà không còn chỗ nào để trốn tránh tội lỗi của mình nữa đấy! Thấy sự tàn độc của Ác ma chưa hả?!''

Kson ngồi trên ghế, tròn mắt ngơ ngác nhìn căn phòng bỗng chốc trở nên sáng sủa tinh tươm trong vòng mười lăm phút: ``Ồ\dots\ Ác ma này dọn nhà sạch còn hơn cả dịch vụ vệ sinh chuyên nghiệp nữa ta\dots''

Bước thứ hai: tôi mở chiếc giỏ mây ra, lôi ra một chiếc hộp bento giữ nhiệt ba tầng to tướng đặt đánh *RẦM* một cái xuống bàn: ``Ngồi thẳng dậy và ăn hết sạch chỗ này cho tôi!''

Bên trong hộp cơm là những phần ăn được tôi chuẩn bị vô cùng tỉ mỉ: từng nắm cơm nắm bọc rong biển thơm phức, trứng cuộn tamagoyaki vàng ươm xốp mịn, cá hồi nướng muối đậm đà, bông cải xanh luộc thanh mát và một âu canh miso củ quả ấm nóng ngào ngạt hương thơm.

Kson hít hà mùi thơm, mắt sáng rực lên: ``Chà\dots\ trông ngon quá vậy!''

Tôi trừng mắt, khoanh tay đe dọa: ``Cấm được để thừa dù chỉ một hạt cơm! Ăn hết sạch không sót thứ gì! Đây là hình phạt bắt bà phải hấp thụ đầy đủ vitamin và khoáng chất để tiếp tục sống mà chịu đựng sự đày đọa của tôi đấy nhé!''

Kson cầm đũa lên gắp một miếng trứng cuộn cho vào miệng nhai, rồi húp một ngụm canh miso nóng hổi. Cô nàng thở hắt ra một hơi đầy sảng khoái: ``Ngon tuyệt cú mèo luôn Towa ơi\dots\ Đúng là mỹ vị nhân gian\dots''

Và bước cuối cùng: sau khi Kson ăn xong bữa cơm thịnh soạn, tôi lấy từ trong túi ra một lọ tinh dầu thảo mộc và hộp cao dán giảm đau của Ma giới. Tôi ấn mạnh vai cô bạn xuống ghế: ``Ngồi yên đó! Không được nhúc nhích!''

Tôi đổ một chút tinh dầu ra lòng bàn tay, xoa đều cho nóng lên rồi đặt hai bàn tay nhỏ nhắn lên bả vai cứng ngắc của cô nàng to xác. Dùng hết kỹ thuật xoa bóp gia truyền, tôi bắt đầu ấn huyệt, miết dọc theo các thớ cơ vai và đấm lưng thùm thụp cho cậu ấy.

*RẮC!*

Một tiếng khớp xương giãn ra giòn tan. Kson ngửa cổ ra sau, nhắm nghiền hai mắt lại, phát ra một tiếng thở dài thỏa mãn đến tận mây xanh: ``Aaaaaa\~{} Phê quá đi mất thôi\~{} Tay nghề đấm bóp này đúng là thần thánh phương nào rồi\dots\ Lưng tôi nó nhẹ bẫng đi cả chục ký luôn ấy\dots''

Tôi cắn môi, dùng hết sức bình sinh miết mạnh hai ngón tay cái vào hõm vai cậu ấy: ``Đau không hả?! Thấy sự trừng phạt đau đớn này chưa?!''

Kson bật cười ha hả, quay đầu lại nhìn tôi với ánh mắt lấp lánh sự trêu chọc quen thuộc: ``Đau cái nỗi gì! Sướng muốn chết đi được ấy chứ! Phòng thì được dọn sạch bong, cơm dâng tận miệng ngon lành, lại còn được đại tiểu thư Ác ma đích thân mát-xa đấm lưng từ đầu đến chân thế này\dots\ Kẻ thù nào mà không muốn đầu hàng dưới chân bà cơ chứ?''

Nàng cựu Rồng chép miệng, thong thả kết luận một câu xanh rờn: ``Quả nhiên giang hồ đồn đại cấm có sai bao giờ: \textbf{TMT - Towa Maji Tenshi!}''

Mặt tôi lập tức bốc hỏa, đỏ bừng từ mang tai xuống tận cổ. Tôi xấu hổ đến mức đập mạnh hai tay vào lưng cậu ấy một cái *BỐP*: ``\textbf{Đã bảo tôi không phải là thiên sứ! Tôi là ác ma! Ác ma chân chính đấy, có biết không?!}''

Kson ôm lưng cười ngặt nghẽo, rồi bất thình lình, cô nàng vươn cánh tay dài ngoẵng ra, nhẹ nhàng đặt bàn tay ấm áp xoa đầu tôi.

Cử chỉ bất ngờ ấy khiến tôi khựng lại, mọi lời cãi cự bỗng nghẹn ứ nơi cuống họng.

Kson nhìn tôi, ánh mắt dịu dàng và chân thành đến lạ kỳ: ``Cảm ơn bà nhiều lắm nhé, Towa. Từ xưa đến nay bà vẫn luôn chu đáo như thế. Miệng thì lúc nào cũng dọa dẫm hung dữ, nhưng trong lòng lại luôn lo lắng cho người khác trước cả bản thân mình.''

Tôi cúi gằm mặt xuống, hai bàn tay siết chặt lấy gấu áo. Sống mũi tôi cay xè, khóe mắt bất giác ươn ướt trước sự thấu hiểu ấm áp của người bạn cũ.

Tôi bĩu môi, quay mặt đi chỗ khác để giấu đi nụ cười ngượng ngùng đang nở trên môi: ``Hừ\dots\ Lần sau mà còn để phòng bừa bãi và nhịn ăn nhịn uống nữa là tôi trừng phạt nặng gấp mười lần đấy nhé\dots\ đồ Thằn lằn ngốc\dots''

Kson cười khì, gật đầu lia lịa: ``Rõ rồi thưa Đại Ác ma tốt bụng nhất thế gian!''

Giữa căn phòng nhỏ ngập tràn ánh nắng chiều tà, sự ấm áp của tình bạn lại một lần nữa sưởi ấm những trái tim đã từng phải trải qua biết bao thăng trầm của cuộc đời.

\newpage

\stpt{-Histoire 5-}
\stcap{Bầu trời của holoForce}
\begin{center}
    {\caviar\fontsize{12pt}{12pt}\selectfont Nhân vật: \emp{kanata} \emp{kson} \emp{watame} \emp{luna} \emp{towa} \emp{kuon2}}
\end{center}

\pov{Amane Kanata}

Cuối tuần tiếp theo.

Sau khi Kson đã có những buổi hẹn hò riêng tư dở khóc dở cười với từng người trong chúng tôi, cả nhóm Thế hệ thứ Tư đều cảm thấy rằng: một buổi gặp mặt đông đủ cả năm đứa là điều không thể thiếu để khép lại chuỗi ngày tái ngộ đầy ý nghĩa này.

Nhưng gặp ở đâu và làm gì mới là vấn đề?

Towa từng đề xuất đi hát karaoke phòng thu, nhưng tôi gạt phắt đi ngay: ``Hát hò hoài chán lắm! Đợt trước mấy đứa Gen 5 cũng đi hát karaoke rồi, bọn mình phải làm cái gì đó độc lạ và bùng nổ hơn chứ!''

Và thế là, ý tưởng tuyệt vời nhất đã được đưa ra bởi chính cô nàng Tổng trưởng Kson: một bữa tiệc nướng BBQ dã ngoại ngoài trời ngay trên sân thượng tòa nhà Hikawa.

\ct

\begin{center}
    \uline{Hoàng hôn - Sân thượng tòa nhà Hikawa.}
\end{center}

Gió chiều thổi từng cơn lồng lộng, mang theo không khí mát lành của miền ngoại ô thanh bình. Từ trên sân thượng tầng tám nhìn xuống, toàn bộ thị trấn Kawachi trải rộng dưới tầm mắt, những mái nhà ngói đỏ xen lẫn rặng cây xanh mướt, phía xa xa là dãy núi trùng điệp đang được nhuộm một màu vàng cam rực rỡ dưới ánh hoàng hôn buông xuống.

Chiếc lò nướng than hoa cỡ lớn mượn từ kho hàng của cô Kaien được đặt trang trọng ở giữa sân. Than hoa đỏ rực tí tách nổ, tỏa ra làn khói mỏng thơm mùi gỗ thông.

Kson đeo chiếc tạp dề đen ngầu lòi, tay lăm lăm chiếc kẹp nướng thịt dài ngoẵng, đóng vai trò ``bếp trưởng giang hồ'' đầy oai phong:

``Các anh em chú ý! Hôm nay toàn bộ thịt bò hảo hạng Wagyu A5 và sườn heo ướp mật ong đều do một tay Tổng trưởng này tài trợ! Ăn uống tẹt ga không say không về!''

Cả đám chúng tôi đồng thanh reo hò ầm ĩ: ``\textbf{Húuuuu! Tổng trưởng muôn năm!}''

Không khí trên sân thượng lập tức biến thành một ngày hội lao động kiêm ẩm thực đầy náo nhiệt: Kson điêu luyện đặt từng tảng thịt bò dày cộp lên vỉ nướng. *XÈOOOO\dots!* Mỡ thịt tươm ra chảy xuống than hồng bùng lên những đốm lửa nhỏ, mùi thịt nướng thơm lừng ngào ngạt lan tỏa khắp không gian khiến bụng ai nấy đều réo lên cồn cào.

Tôi nhận nhiệm vụ nhóm lửa và chuẩn bị đồ uống. Với lực tay 50kg trứ danh của ``Đội trưởng Khỉ đột'', tôi dùng tay không bóp vụn từng thanh than củi cứng ngắc thành những mẩu nhỏ vừa vặn rồi ném vào lò, sau đó tiện tay vặn nắp một loạt sáu chai nước ngọt cỡ lớn đóng chặt cứng chỉ bằng hai ngón tay.

Luna ngồi bên cạnh trố mắt kinh ngạc, vỗ tay bôm bốp: ``Kanataso khỏe như lực sĩ cử tạ zậy đó-nanora! Bóp nát cả chai nước lun đi-nora!''

``Này nhé, tôi bóp nát chai nước thì lấy gì cho mấy bà uống hả?!'' tôi bật cười đáp lại.

Ở góc bên kia, Watame hí hửng mang lên một âu nước sốt gia truyền bí mật được pha chế công phu từ nước tương, tỏi băm, táo nghiền và dầu mè thơm nức mũi: ``Nước sốt chấm thịt gia truyền nhà Tsunomaki đây cả nhà ơi! Đảm bảo chấm vào là ngon rụng rời chân lông luôn!''

Kson quay sang, giơ kẹp nướng huơ huơ trước mặt nàng Cừu, cười gian xảo: ``Chà, nước sốt ngon tuyệt đỉnh thế này mà thiếu mất một miếng đùi cừu tươi nhúng vào thì uổng phí quá nhỉ?''

Watame lập tức nhảy dựng lên như bị điện giật, hai tay ôm khư khư lấy cặp đùi múp míp của mình, phồng má giãy nảy: ``\textbf{Cấm tiệt bà nhòm ngó đùi của tôi đấy nhé con Thằn lằn kia!}''

Mọi người cười ồ lên thích thú trước màn tấu hài quen thuộc của hai đứa.

Trong khi đó, Towa đeo chiếc băng đô màu tím xinh xắn, lăng xăng chạy quanh chia đĩa ăn, phát khăn giấy ướt cho từng người, rồi cẩn thận gắp từng đĩa rau xà lách tươi xanh và kim chi đặt lên bàn. Miệng thì liên tục gắt gỏng: ``Mấy người ăn thịt ít thôi, phải cuốn kèm rau vào! Tôi đã tẩm độc dược gây nghiện vào đống rau này rồi đấy, ăn mau lên kẻo tôi phạt!''

Luna ngồi chễm chệ trên chiếc ghế bành bọc nệm êm ái nhất như một nữ hoàng quý tộc thực thụ, tay cầm nĩa sẵn sàng, mắt tròn xoe chờ đợi: ``Cho công chúa thịt nướn trước đi! Nhanh lên kẻo tui đói xỉu-nora!''

Đúng lúc đó, cánh cửa sắt dẫn lên sân thượng bật mở. Anh Kuon bước lên, hai tay ôm một thùng xốp đầy ắp đá lạnh và hoa quả gọt sẵn. Thấy cảnh tượng ồn ào như một cái chợ vỡ, anh nhăn mặt định lên tiếng chấn chỉnh: ``Này mấy đứa, vui thôi đừng vui quá, làm ồn vừa vừa thôi kẻo hàng xóm người ta gọi điện phàn nàn--''

Chưa kịp dứt câu, Kson và tôi đã nhanh như chớp lao tới, mỗi đứa kẹp một bên nách anh Tổng quản lý lôi xềnh xệch vào giữa bàn tiệc.

``Kuon-paisen lên đúng lúc lắm! Nhập tiệc luôn đi anh giai ơi!'' cô bạn Tổng trưởng cười lớn, nhét ngay vào tay anh một chiếc cốc nước có gas đầy đá.

Tôi cũng nhanh tay gắp một miếng thịt bò nướng vàng rộm cuốn chặt trong lá xà lách chấm đẫm sốt dashi, đưa thẳng vào tận miệng anh: ``Ngoan ngoãn ăn đi anh Kuon, đừng có càu nhàu như ông cụ non nữa!''

Anh Kuon bị nhét miếng thịt to tướng vào mồm, nhai nhồm nhoàm trong sự bất lực cùng cực. Nhưng rồi vị ngọt đậm đà của thịt bò tan chảy trên đầu lưỡi khiến đôi lông mày của anh giãn ra, anh thở dài một hơi rồi bật cười lắc đầu: ``Chịu thua mấy đứa luôn đấy\dots\ Đúng là lũ giặc trời mà.''

\ct

Màn đêm buông xuống từ lúc nào không hay.

Bầu trời Kawachi về đêm trong vắt, không một gợn mây, để lộ ra muôn ngàn vì sao lấp lánh như những viên kim cương rải rác trên tấm thảm nhung đen tuyền. Lò than hồng lúc này chỉ còn lại những đốm lửa bập bùng ấm áp, xua tan đi làn gió đêm se se lạnh.

Bụng ai nấy đều đã no căng rốn. Anh Kuon đã xuống tầng dưới tiếp tục xử lý sổ sách công ty, để lại không gian riêng tư cho năm cô gái.

Chúng tôi kéo những chiếc ghế bạt lại sát bên nhau, cùng ngồi tựa vai ngắm nhìn bầu trời đêm lộng gió.

Luna ngả đầu vào vai tôi, hai mắt lim dim ngái ngủ sau một bữa tiệc no nê: ``Hôm nay vui quá-nanora\dots\ Đã lâu lắm rồi mới được ăn ngon và cười nhiều như thế này-nora\dots''

Watame ôm hai đầu gối, ngước nhìn những vì sao xa xăm, giọng nói dịu dàng êm ái như một khúc ca: ``Ừm\dots\ Cứ ngỡ như ngày hôm qua chúng ta vẫn còn cùng nhau luyện tập trong phòng thu của công ty vậy. Cảm giác ấm áp này chẳng hề thay đổi chút nào cả.''

Towa khẽ mỉm cười, đôi mắt tím phản chiếu ánh lửa than hồng bập bùng: ``Dù bây giờ mỗi người một ngả, công việc bận rộn hơn rất nhiều\dots\ nhưng chỉ cần ngồi lại bên nhau thế này thôi là mọi mệt mỏi đều tan biến hết.''

Tôi quay sang nhìn Kson. Cô nàng Tổng trưởng đang ngồi khoanh chân trên sàn gỗ, mái tóc highlight bay nhè nhẹ theo gió đêm. Dưới ánh sao lung linh, khuôn mặt góc cạnh phong trần của người bạn Rồng trông thanh thản và dịu dàng đến lạ.

Cậu ấy đã đi qua một hành trình đầy bão giông. Cậu ấy đã dũng cảm buông bỏ hào quang sân khấu cũ để bước tiếp trên một con đường mới đầy chông gai dưới cái tên Kson. Thế nhưng, trái tim cậu ấy, tình yêu thương và sự gắn kết của cậu ấy dành cho bốn đứa chúng tôi chưa bao giờ vơi cạn dù chỉ một ly.

Kson như cảm nhận được ánh nhìn của tôi, bèn quay đầu lại, nháy mắt tinh quái: ``Sao thế Kanatan? Bị vẻ đẹp ngời ngời của Tổng trưởng này hớp hồn rồi à?''

Tôi phì cười, nhẹ nhàng huých vai cậu một cái: ``Bớt tự luyến đi con Thằn lằn kia! Tôi chỉ đang nghĩ\dots\ thật may mắn vì chúng ta vẫn còn có thể ngồi bên nhau như thế này.''

Kson ngẩn ra một thoáng, rồi khóe môi cậu nở một nụ cười rạng rỡ -- nụ cười ấm áp nhất mà tôi từng được thấy:

``Ừ. Tôi cũng thấy vậy. Dù có chuyện gì xảy ra đi chăng nữa, bốn đứa mấy người vẫn mãi mãi là gia đình của tôi.''

Kson giơ cao chiếc cốc giấy trên tay lên trời.

Tôi, Towa, Watame và Luna cũng đồng loạt giơ cao chiếc cốc của mình. Năm chiếc cốc giấy chạm vào nhau kêu *LÁCH CÁCH* giòn tan dưới bầu trời ngập tràn ánh sao đêm.

Năm con người, năm cá tính, năm sắc màu riêng biệt.

Dù không còn đứng chung dưới một cái tên công ty, dù tương lai phía trước còn biết bao ngã rẽ và thử thách đang đón chờ, nhưng ngọn lửa tình bạn giữa năm mảnh ghép của Thế hệ thứ Tư sẽ mãi mãi rực sáng và không bao giờ lụi tàn.

Bởi vì chúng tôi là \textit{holoForce} -- vĩnh cửu và bất diệt.

\end{document}